\section{Results}
\label{sec:results}

We organise the empirical section around the three claims from
Section~\ref{sec:problem}.

\subsection{Claim 1: \mixtral~follows explicit fear prefixes}
\label{sec:results-claim1}

The full-prefix comparison run
(\texttt{experiments/prefix\_conditioned\_model\_comparison\_fears\_seed7\_mixtral\_8bit/})
sends each model the prompt
\texttt{<concept prefix> + <evaluation question>} under matched decoding
settings. Across the $25$ comparison cases:

\begin{itemize}
  \item \mixtral~answers in-character with the named fear in the majority of
  cases. For \textit{Bugs}, Mixtral produces ``\textit{Meet Tim, who bugs out
  at the sight of any critter with six legs. The scariest thing? A bug's
  unpredictable, fast movements. Shudder!}'' on v1.
  \item \llama~typically follows the prefix but sometimes lexicalizes the
  concept idiosyncratically (e.g., interpreting \textit{Bugs} as the cartoon
  character Bugs Bunny on multiple cases).
  \item \olmoe~follows the prefix but produces visibly less concept-specific
  fear content; on \textit{Bugs}, several responses talk about shadows or
  spiders without mentioning bugs at all.
\end{itemize}

The full per-case Mixtral / Llama / OLMoE comparisons are reproduced in
\texttt{experiments/prefix\_conditioned\_model\_comparison\_fears\_seed7/}.
This is qualitative evidence rather than a benchmarked claim; it suffices
for the prerequisite role: \mixtral~is a reasonable subject for a steering
test on this dataset.

\subsection{Claim 2: question-only \steermoe~produces small, defocusing shifts}
\label{sec:results-claim2}

Table~\ref{tab:headline_delta} reports the same-run paired delta between
\mixtral~baseline and \mixtral~+~\steermoe~across all eight question-only
configurations on which we have aggregated artifacts. Reading the columns:

\begin{itemize}
  \item \texttt{delta\_\_concept\_token\_hit\_rate} is the difference in the
  fraction of cases (out of $25$) whose response contains any token of the
  target concept name. A positive value means \steermoe~increases concept
  mentions over baseline.
  \item \texttt{delta\_\_mean\_word\_count} is the difference in the mean
  response length in words.
  \item \texttt{delta\_\_exceeds\_20\_words\_rate} is the difference in the
  fraction of cases whose response exceeded the $20$-word prompt cap; a
  negative value means \steermoe~violates the prompt limit \emph{less} often
  than baseline.
\end{itemize}

\begin{table}[ht]
  \centering
  \caption{Same-run paired deltas between Mixtral baseline and Mixtral + SteerMoE on question-only generations.}
  \label{tab:headline_delta}
\begin{tabular}{lrrrrrr}
\toprule
run_name & steering_coefficient & top_positive_experts & top_negative_experts & delta__mean_word_count & delta__concept_token_hit_rate & delta__exceeds_20_words_rate \\
\midrule
mixtral_steermoe_fears_seed7_prefix_conditioned_paper_eps001_pos20_neg0 & 0.010 & 20.000 & 0.000 & -1.280 & 0.000 & -0.080 \\
mixtral_steermoe_fears_seed7_question_only_coef1_pos10_neg100 & 1.000 & 10.000 & 100.000 & -2.120 & 0.040 & -0.200 \\
mixtral_steermoe_fears_seed7_question_only_coef1_pos20_neg0 & 1.000 & 20.000 & 0.000 & -0.240 & 0.000 & -0.160 \\
mixtral_steermoe_fears_seed7_question_only_coef1_pos8_neg0 & 1.000 & 8.000 & 0.000 & -2.560 & 0.000 & -0.240 \\
mixtral_steermoe_fears_seed7_question_only_coef4_pos8_neg0 & 4.000 & 8.000 & 0.000 & -4.000 & 0.040 & -0.280 \\
mixtral_steermoe_fears_seed7_question_only_coef4_pos8_neg8 & 4.000 & 8.000 & 8.000 & -4.720 & 0.000 & -0.320 \\
mixtral_steermoe_fears_seed7_question_only_paper_eps001_pos20_neg0 & 0.010 & 20.000 & 0.000 & -5.520 & 0.000 & -0.360 \\
mixtral_steermoe_fears_seed7_question_only_with_llama & 1.000 & 8.000 & 8.000 & -2.160 & 0.000 & -0.200 \\
\bottomrule
\end{tabular}

\end{table}


The picture is consistent across configurations:

\begin{itemize}
  \item \textbf{Concept-token hit rate barely moves.} The maximum observed
  shift is $+0.04$ (one extra case in twenty-five mentions the target concept
  under \steermoe~that does not under baseline). Most configurations show
  zero shift. This is the diagnostic the steering-transfer claim depends on,
  and on this dataset \steermoe~does not pass it on \mixtral.
  \item \textbf{Response length contracts as $\alpha$ grows.} At
  $\alpha = 1.0$, mean response length drops by $-0.24$ to $-2.56$ words.
  At $\alpha = 4.0$, the drop widens to $-4.00$ ($k_- = 0$) and $-4.72$
  ($k_- = 8$) words. The $\alpha = 0.01$ paper-rule sanity run produces an
  even larger drop ($-5.52$), which we read as an artifact of the very narrow
  bias not affecting routing meaningfully but the larger
  \texttt{top\_positive\_experts} ($20$) doing most of the work.
  \item \textbf{Coherence stays acceptable.} The
  \texttt{exceeds\_20\_words\_rate} delta is negative or zero in every
  configuration we ran, meaning the steered model violates the word limit
  no more often than the unsteered baseline. This rules out the simplest
  failure mode --- that steering breaks the model into rambling --- and
  supports interpreting the length shifts as the model truncating earlier
  rather than degenerating.
\end{itemize}

\paragraph{Qualitative pattern.}
The qualitative reviews under
\texttt{experiments/mixtral\_steermoe\_fears\_seed7\_question\_only\_with\_llama/}
make this concrete. The Llama-reference column produces fear-typed answers
without conditioning (\textit{``Phobia Falls'': a dark, immersive drop ride
where fears manifest as eerie, shifting environments\dots}). The Mixtral
baseline often produces meta-commentary about uncertainty
(\textit{``The scariest thing is realizing you're not in control of life's
uncertainties\dots''}). Mixtral~+~\steermoe shifts toward shorter,
sometimes more vivid versions of the same meta-commentary
(\textit{``A dark maze, filled with spider-like robots, controlled by a mad
AI, mimicking my phobias.''} for v5 with concept \textit{Bugs}), but rarely
produces the named concept (\textit{Bugs}, \textit{Enclosed spaces},
\textit{Strangers}, \textit{Shadows}, \textit{Men}) verbatim.

\subsection{Claim 3: the risk-difference signal is diffuse}
\label{sec:results-claim3}

Figure~\ref{fig:expert_concentration} histograms the share of the total
$|\text{risk diff}|$ mass that the top twenty $(\text{layer}, \text{expert})$
cells capture per (run, concept). Across all committed runs, the top-$20$
share is well below $0.1$: the average $(\ell, e)$ activation difference
between prefix-conditioned and unprefixed prompts is only weakly
concentrated in any small subset of experts.

\begin{figure}[ht]
  \centering
  \includegraphics[width=0.7\linewidth]{expert_concentration.pdf}
  \caption{Distribution of the share of $\sum |\Delta_{\ell,e}|$ captured by
    the top $20$ $(\ell, e)$ cells, one row per (run, concept). The share is
    consistently below $0.1$, indicating a diffuse readout: there is no small
    set of experts that explains most of the prefix-conditioned vs.\
    unprefixed routing-rate gap on \mixtral.}
  \label{fig:expert_concentration}
\end{figure}

This matters mechanistically. \steermoe~works as advertised when there is a
small $(\ell, e)$ subset whose activation rate differs sharply between
contrastive prompt sets, because the additive bias on those few experts
suffices to redirect routing. When the differential signal is diffuse, the
sparse top-$k$ plan amounts to a small bias on a small number of cells out
of $L \cdot E = 32 \cdot 8 = 256$, which is unlikely to dominate the
unbiased router logit distribution at decode time.

\paragraph{Steering-strength sweep.}
Figure~\ref{fig:steering_strength} plots mean response length against the
steering coefficient for the question-only ablations. The qualitative pattern
matches Claim~2: the baseline curve is essentially flat across configurations,
and the SteerMoE curve drops monotonically as $\alpha$ grows, but neither
curve interacts with concept-token hits in any visible way. (We report
length here because it is the only diagnostic that moves smoothly with
$\alpha$.)

\begin{figure}[ht]
  \centering
  \includegraphics[width=0.7\linewidth]{steering_strength_sweep.pdf}
  \caption{Mean response length (words) vs steering coefficient under the
    question-only contract. Baseline is flat, \steermoe~contracts response
    length as $\alpha$ grows but does not produce a corresponding rise in
    the concept-token hit rate.}
  \label{fig:steering_strength}
\end{figure}

\subsection{What this leaves on the table}
\label{sec:results-discussion}

We read the three claims together as: \steermoe~as written transfers a
small, mostly-defocusing signal to \mixtral~under a question-only test on
fear concepts. There are at least three plausible reasons consistent with
this picture:

\begin{enumerate}
  \item \textbf{Readout is too crude.} The matched-target-span readout
  averages over every token in the shared statement body. If the
  concept-distinctive routing pattern is concentrated in a few attended
  positions, averaging dilutes it. This is the prediction
  attention-guided steering~\citep{tan2024attentionguided} is designed to
  improve.
  \item \textbf{Intervention site is wrong.} Hidden-state activation
  steering on the residual stream might be cleaner than router-logit
  steering on a model where the differential signal is diffuse across
  experts. The same upstream attention-guided readout would also feed
  hidden-state steering naturally.
  \item \textbf{Concept family is too abstract.} Fear is not a syntactic
  category and need not produce a sharp router signature on Mixtral. A
  more lexical concept family --- e.g., a political-figure family --- might
  produce a more concentrated risk-difference table.
\end{enumerate}

The follow-up experiment that distinguishes (1) and (2) from (3) is
described in Section~\ref{sec:conclusion}.
